% ---------- 扉页（与原书一致） ----------
\begin{titlepage}
  \centering
  \vspace*{1.5cm}
  {\large Lecture Notes in Mathematics\par}
  \vspace{0.6em}
  {\small An informal series of special lectures, seminars and reports on mathematical topics\par}
  \vspace{1.2em}
  {\small Edited by A.\ Dold, Heidelberg and B.\ Eckmann, Z\"urich\par}
  \vspace{3em}
  {\Large 21\par}
  \vspace{4em}
  {\normalsize A.\ Borel \textperiodcentered\ S.\ Chowla \textperiodcentered\ C.\ S.\ Herz\\
   K.\ Iwasawa \textperiodcentered\ J-P.\ Serre\par}
  \vspace{3em}
  {\huge Seminar\\[0.35em] on Complex Multiplication\par}
  \vspace{2.5em}
  {\small Seminar held at the Institute for Advanced Study, Princeton, N.J., 1957--58\par}
  \vspace{2em}
  {\normalsize 1966\par}
  \vfill
  Springer-Verlag \textperiodcentered\ Berlin \textperiodcentered\ Heidelberg \textperiodcentered\ New York
  \vspace{1.2em}

  {\large 译者\quad 翻译及重排\par}
\end{titlepage}

% 原书内封前的丛书说明（Lecture Notes 丛书宗旨，英文部分）
\noindent\frenchspacing
This series aims at speedy, informal, and high level information on new developments in mathematical research and teaching. Considered for publication are:
\begin{enumerate}[leftmargin=2.5em]
  \item Preliminary drafts of original papers and monographs
  \item Special lectures on a new field, or a classical field from a new point of view
  \item Seminar reports
  \item Reports from meetings
\end{enumerate}
Out of print manuscripts satisfying the above characterization may also be considered, if they continue to be in demand.

\medskip
The timeliness of a manuscript is more important than its form, which may be unfinished and preliminary. In certain instances, therefore, proofs may only be outlined, or results may be presented which have been or will also be published elsewhere.

\medskip
The publication of the ``Lecture Notes'' Series is intended as a service, in that a commercial publisher, Springer-Verlag, makes house publications of mathematical institutes available to mathematicians on an international scale. By advertising them in scientific journals, listing them in catalogs, further by copyrighting and by sending out review copies, an adequate documentation in scientific libraries is made possible.

\section*{Manuscripts}
Since manuscripts will be reproduced photomechanically, they must be written in clean typewriting. Handwritten formulae are to be filled in with indelible black or red ink. Any corrections should be typed on a separate sheet in the same size and spacing as the manuscript. All corresponding numerals in the text and on the correction sheet should be marked in pencil. Springer-Verlag will then take care of inserting the corrections in their proper places. Should a manuscript or parts thereof have to be retyped, an appropriate indemnification will be paid to the author upon publication of his volume. The authors receive $25$ free copies.

\medskip
Manuscripts in English, German or French should be sent to Prof.\ Dr.\ A.\ Dold, Mathematisches Institut der Universit\"at Heidelberg, Tiergartenstra\ss{}e or Prof.\ Dr.\ B.\ Eckmann, Eidgen\"ossische Technische Hochschule, Z\"urich.

\clearpage

% ==================== 目录 ====================
\tableofcontents
